\chapter{Profilo Utente e Persona}

\section{Profilo utente}

L'applicativo è progettato per due tipologie di utente:
\begin{itemize}
    \item \textbf{Configuratore}: rappresenta l'attore con ruolo di amministratore e gestore del sistema, non prende parte attiva agli scambi, ma ha il compito di predisporre, governare e supervisionare l'ambiente operativo.
    \item \textbf{Fruitore}: è l'attore che si interfaccia con l'applicazione al fine di proporre e finalizzare scambi con altri fruitori.
\end{itemize}

\subsection{Configuratore}
\subsubsection{Caratteristiche psicologiche}

\begin{table}[H]
    \centering
    \begin{tabular}{@{}>{\raggedright\arraybackslash}p{4cm}>{\raggedright\arraybackslash}p{8cm}@{}}
        \toprule
        \textbf{Caratteristica} & \textbf{Tipologia}                        \\ \midrule
        Stile cognitivo         & Astratto, Deduttivo, Spaziale e Analitico \\
        Attitudine              & Positiva / Neutrale                       \\
        Motivazione             & Alta                                      \\ \bottomrule
    \end{tabular}
    \caption{Caratteristiche psicologiche del configuratore}
\end{table}
Il configuratore non opera su oggetti fisici individuali, bensì su modelli e schemi logici, applicando un ragionamento top-down. Deve concettualizzare insiemi, gerarchie ad albero (categorie radice e sottocategorie) e attributi (campi obbligatori vs. facoltativi) per questo presenta uno stile cognitivo \emph{astratto} e \emph{deduttivo}.
Dato il ruolo amministrativo di verificare la coerenza dei dati deve compiere decisioni \emph{analitiche}, non intuitive, poiché un errore di configurazione si ripercuote su tutti gli utenti.

Il sistema è per il configuratore uno strumento di lavoro, questo implica che la sua attitudine sia positiva o al più neutrale rispetto al sistema. Per questo motivo è altamente motivato all'uso dell'applicazione.

\subsubsection{Caratteristiche di conoscenza}
\begin{table}[H]
    \centering
    \begin{tabular}{@{}>{\raggedright\arraybackslash}p{8cm}>{\raggedright\arraybackslash}p{4cm}@{}}
        \toprule
        \textbf{Caratteristica}      & \textbf{Tipologia} \\ \midrule
        Livello di alfabetizzazione  & Buon livello       \\
        Titolo di studio             & Superiore o Laurea \\
        Abilità dattilografica       & Medio-alta         \\
        Linguaggio nativo            & Italiano           \\
        Alfabetizzazione informatica & Alta               \\ \bottomrule
    \end{tabular}
    \caption{Caratteristiche di conoscenza del configuratore}
\end{table}

Il configuratore presenta un \emph{buon livello di alfabetizzazione} e un titolo di studio tipicamente pari a \emph{scuola superiore o laurea}. Anche il livello di \emph{alfabetizzazione informatica è alto}. Questa formazione è necessaria per comprendere e modellare l'architettura logica del catalogo.

La sua abilità dattilografica è \emph{medio-alta}, poiché l'attività di amministrazione richiede l'inserimento frequente di dati testuali.

Il linguaggio nativo di riferimento è l'\emph{italiano}, lingua in cui è formulata la terminologia della piattaforma.

\subsubsection{Caratteristiche di esperienza}
\begin{table}[H]
    \centering
    \begin{tabular}{@{}>{\raggedright\arraybackslash}p{8cm}>{\raggedright\arraybackslash}p{4cm}@{}}
        \toprule
        \textbf{Caratteristica}            & \textbf{Tipologia} \\ \midrule
        Esperienza del compito             & Medio/Esperto      \\
        Esperienza sull'applicazione       & Medio/Esperto      \\
        Esperienza dei sistemi informatici & Alta               \\
        Uso di altri sistemi               & Frequente          \\ \bottomrule
    \end{tabular}
    \caption{Caratteristiche di esperienza del configuratore}
\end{table}

Il configuratore possiede un'\emph{esperienza del compito da esperto}: ha una solida padronanza delle regole organizzative che governano il sistema, sa come strutturare un catalogo.

Per quanto concerne l'\emph{esperienza sull'applicazione}, si colloca tipicamente a livello \emph{medio o alto}. Grazie alla comprensione del dominio, raggiunge rapidamente un livello medio/alto di padronanza dell'interfaccia.

L'\emph{esperienza dei sistemi informatici è alta come la loro frequenza di utilizzo} : il configuratore è abituato a interagire con computer, file system, finestre modali, software di produttività e formati strutturati (es. JSON), comprendendo senza difficoltà messaggi di sistema.

\subsubsection{Caratteristiche fisiche}
\begin{table}[H]
    \centering
    \begin{tabular}{@{}>{\raggedright\arraybackslash}p{4cm}>{\raggedright\arraybackslash}p{8cm}@{}}
        \toprule
        \textbf{Caratteristica} & \textbf{Tipologia}              \\ \midrule
        Daltonismo              & Sì / No                         \\
        Manualità               & Destrorso, Mancino o Ambidestro \\
        Disabilità              & Assenti o lievi                 \\ \bottomrule
    \end{tabular}
    \caption{Caratteristiche fisiche del configuratore}
\end{table}

Il profilo fisico del configuratore è quello tipico di un operatore d'ufficio o amministrativo su postazione desktop.

\subsubsection{Caratteristiche del lavoro e dei compiti}
\begin{table}[H]
    \centering
    \begin{tabular}{@{}>{\raggedright\arraybackslash}p{6cm}>{\raggedright\arraybackslash}p{6cm}@{}}
        \toprule
        \textbf{Caratteristica}          & \textbf{Tipologia}                       \\ \midrule
        Uso del sistema                  & Obbligatorio                             \\
        Frequenza d'uso                  & Bassa / Media                            \\
        Addestramento di base            & Sui manuali / Formale                    \\
        Categorie di lavoro              & Tecnico o Impiegato    \\
        Uso di altri strumenti           & Altro \\
        Frequenza del turn-over          & Basso                                    \\
        Importanza del compito           & Alta                                     \\
        Struttura del compito            & Alta                                     \\ \bottomrule
    \end{tabular}
    \caption{Caratteristiche del lavoro e dei compiti del configuratore}
\end{table}

L'\emph{uso del sistema è obbligatorio} perché la configurazione dell'applicazione rientra tra i doveri assegnati al ruolo. La \emph{frequenza d'uso è bassa o media} in quanto l'interazione è concentrata nella fase di avvio e popolamento iniziale del catalogo.

Data la delicatezza delle impostazioni e importanza del compito il configuratore fa affidamento per l'addestramento alla documentazione tecnica.

\subsection{Fruitore}
\subsubsection{Caratteristiche psicologiche}
\begin{table}[H]
    \centering
    \begin{tabular}{@{}>{\raggedright\arraybackslash}p{4cm}>{\raggedright\arraybackslash}p{8cm}@{}}
        \toprule
        \textbf{Caratteristica} & \textbf{Tipologia}                       \\ \midrule
        Stile cognitivo         & Concreto, Induttivo, Verbale e Intuitivo \\
        Attitudine              & Positiva                                 \\
        Motivazione             & Moderata / Alta                          \\ \bottomrule
    \end{tabular}
    \caption{Caratteristiche psicologiche del fruitore}
\end{table}
Il fruitore interagisce con singoli annunci e appuntamenti reali, apprendendo per tentativi ed esempi osservati nella piattaforma; per questo presenta uno stile cognitivo \emph{concreto} e \emph{induttivo}.
Dato l'obiettivo di individuare e scambiare beni di proprio interesse in modo semplice e rapido, compie decisioni prevalentemente \emph{intuitive} e guidate da elementi descrittivi e testuali (\emph{verbale}), basandosi su impressioni immediate e a basso carico cognitivo piuttosto che su un'analisi formale dell'architettura sottostante (come invece fa il configuratore).

Il sistema è per il fruitore un mezzo per ottenere un beneficio personale attraverso il baratto, il che si traduce in un'attitudine tendenzialmente positiva rispetto alla piattaforma. La sua motivazione all'uso è moderata o alta: a differenza del configuratore, è fortemente legata al raggiungimento del proprio scopo, risultando più sensibile alla facilità d'uso e al successo effettivo delle proposte di scambio.

\subsubsection{Caratteristiche di conoscenza}
\begin{table}[H]
    \centering
    \begin{tabular}{@{}>{\raggedright\arraybackslash}p{6cm}>{\raggedright\arraybackslash}p{6cm}@{}}
        \toprule
        \textbf{Caratteristica}      & \textbf{Tipologia}             \\ \midrule
        Livello di alfabetizzazione  & Livello medio (o Buon livello) \\
        Titolo di studio             & Media, Superiore o Laurea      \\
        Abilità dattilografica       & Bassa / Media                  \\
        Linguaggio nativo            & Italiano (o Altro)             \\
        Alfabetizzazione informatica & Media / Bassa                  \\ \bottomrule
    \end{tabular}
    \caption{Caratteristiche di conoscenza del fruitore}
\end{table}

Il profilo del fruitore è molto più eterogeneo rispetto a quello del configuratore. Presenta un \emph{livello medio di alfabetizzazione} e un titolo di studio che spazia \emph{dal diploma di scuola media alla scuola superiore o laurea}. L'interfaccia deve pertanto esprimersi con un linguaggio chiaro, diretto e privo di tecnicismi, garantendo che chiunque sia in grado di leggere i dettagli di un'offerta, comprendere i campi richiesti per inserire un articolo e verificare le condizioni dell'incontro.

L'abilità dattilografica è \emph{bassa o media}: le interazioni testuali dell'utente sono ridotte alla compilazione delle caratteristiche essenziali dell'oggetto (marca, note, stato d'uso).

Il linguaggio nativo prevalente è l'\emph{italiano}, ma il bacino potenziale include utenti di \emph{altra} madrelingua.

L'\emph{alfabetizzazione informatica è media o bassa}: il fruitore ha familiarità con le più comuni applicazioni web e smartphone (social network, app di messaggistica o di compravendita), ma in genere non possiede conoscenze di logica informatica o strutture dati.

\subsubsection{Caratteristiche di esperienza}
\begin{table}[H]
    \centering
    \begin{tabular}{@{}>{\raggedright\arraybackslash}p{7cm}>{\raggedright\arraybackslash}p{5cm}@{}}
        \toprule
        \textbf{Caratteristica}            & \textbf{Tipologia}      \\ \midrule
        Esperienza del compito             & Nuovo nel campo / Medio \\
        Esperienza sull'applicazione       & Novizio                 \\
        Esperienza dei sistemi informatici & Media / Bassa           \\
        Uso di altri sistemi               & Scarso/Frequente        \\ \bottomrule
    \end{tabular}
    \caption{Caratteristiche di esperienza del fruitore}
\end{table}

L'\emph{esperienza del compito del fruitore è eterogenea, spaziando da nuovo nel campo a medio o alta}: la maggior parte degli utenti ha un'idea intuitiva di baratto o ha familiarità con mercatini dell'usato e piattaforme di compravendita tra privati (es. Vinted, Subito), ma molti potrebbero non aver mai gestito operazioni di baratto come quelle offerte da questa applicazione.

Rivolgendosi anche a \emph{novizi} il sistema deve essere progettato con un'elevata facilità di apprendimento secondo il paradigma del "walk-up-and-use", consentendo all'utente di portare a termine lo scambio senza bisogno di consultare manuali o seguire una formazione specifica.

\subsubsection{Caratteristiche fisiche}
\begin{table}[H]
    \centering
    \begin{tabular}{@{}>{\raggedright\arraybackslash}p{4cm}>{\raggedright\arraybackslash}p{8cm}@{}}
        \toprule
        \textbf{Caratteristica} & \textbf{Tipologia}              \\ \midrule
        Daltonismo              & Sì / No                         \\
        Manualità               & Destrorso, Mancino o Ambidestro \\
        Disabilità              & Potenzialmente presenti         \\ \bottomrule
    \end{tabular}
    \caption{Caratteristiche fisiche del fruitore}
\end{table}

Il fruitore rappresenta l'utente generico della popolazione; il suo profilo fisico impone requisiti di accessibilità e usabilità inclusivi.

\subsubsection{Caratteristiche del lavoro e dei compiti}
\begin{table}[H]
    \centering
    \begin{tabular}{@{}>{\raggedright\arraybackslash}p{6cm}>{\raggedright\arraybackslash}p{6cm}@{}}
        \toprule
        \textbf{Caratteristica}          & \textbf{Tipologia}                       \\ \midrule
        Uso del sistema                  & Facoltativo                              \\
        Frequenza d'uso                  & Bassa / Media            \\
        Addestramento di base            & Nessuno                                  \\
        Categorie di lavoro              & Eterogenea  \\
        Uso di altri strumenti           & Altro \\
        Frequenza del turn-over          & Alto                                     \\
        Importanza del compito           & Bassa                \\
        Struttura del compito            & Media                                    \\ \bottomrule
    \end{tabular}
    \caption{Caratteristiche del lavoro e dei compiti del fruitore}
\end{table}

L'\emph{uso del sistema è facoltativo e spesso poco frequente}, infatti, il fruitore accede alla piattaforma su base esclusivamente volontaria, guidato dall'interesse personale a cedere o reperire oggetti mediante il baratto. Per questo motivo \emph{la frequenza di turn-over è alta}.

La \emph{categoria di lavoro è del tutto eterogenea}: il fruitore può essere uno studente, un impiegato, un operaio, un pensionato o un professionista.

\section{Persona}
